\chapter{The inheritance}
\label{ch:inheritance}

\section{A country that was not poor, and was not developed}

On the last day of 1958 Cuba had about 6.6 million people, a fifth of the
population of Spain on an island slightly larger than Portugal, and it was by
most aggregate measures one of the four or five richest countries in Latin
America.\unv{} It had more television sets per head than any country in the
region and among the highest counts in the world; it had one of the densest
railway networks in the hemisphere, inherited from sugar; it had roughly one
physician for every thousand inhabitants, a ratio comparable with much of
Western Europe at the time; it had a literate majority, 76 per cent of the
population over ten by the 1953 census; and it had a life expectancy at birth
somewhere around 62 years \citep{perezjr2015}, higher than every mainland Latin American country
except Argentina and Uruguay.\unv{}

Every one of those numbers is a national average, and every one of them is
concealing the thing that actually mattered, which is that Cuba in 1958 was not
one country but two economies bolted together at the port of Havana.

The urban economy, and above all the capital, was genuinely modern. Havana had
department stores, an air-conditioned bus system, a stock exchange, medical
specialists, a professional press of a dozen daily papers, a university with a
real research tradition, and a middle class of professionals, shopkeepers and
unionised industrial workers whose standard of living would not have been out
of place in southern Europe. Cuban organised labour was strong --- the sugar and
transport unions in particular --- and had extracted job protections and wage
floors that were unusual in the region and that would, a decade later, become
one of the constraints the revolution had to negotiate with rather than around.

The rural economy was something else. The most-quoted single source on it is a
survey of a thousand rural wage-worker families conducted in 1956--57 by the
Agrupaci\'on Cat\'olica Universitaria, a lay Catholic organisation with no
sympathy for what came after, which is why the figures survived scrutiny. It
found that among rural workers 4 per cent ate meat with any regularity, 1 per
cent ate fish, 11 per cent drank milk; that 43 per cent could not read; that
about 60 per cent lived in dwellings with earth floors and thatch roofs; that
under 10 per cent had electricity and about 2 per cent running water \citep{perezstable2012}; and that
the median family's annual income was a fraction of the urban wage.\unv{} The
survey was published in a Catholic magazine in 1957 and its findings were not
seriously contested at the time. Approximately 43 per cent of Cubans lived in
the countryside.

So the honest description of pre-revolutionary Cuba is neither the plundered
colony of the official histories nor the prosperous republic of the exile
memoirs. It is a middle-income country with a first-world capital, a
third-world interior, and a distribution so wide that the national average
describes almost nobody.

\section{Sugar, and the shape it gave everything}

The reason the two economies existed is sugar, and the reason sugar dominated
is a policy instrument in Washington called the quota.

Sugar was roughly four-fifths of Cuban export earnings and had been for
half a century.\unv{} Under the reciprocal trade arrangements that governed
Cuban--American commerce, the United States bought a guaranteed tonnage at a
price above the world market and in return Cuban tariffs favoured American
manufactures. The arrangement was, in the short run, extremely good for Cuba:
it converted a volatile commodity into a stable annuity and it underwrote the
consumption of the urban middle class. In the long run it did three things that
compounded.

It suppressed industry. Why build a Cuban shoe factory when the tariff schedule
made American shoes cheap and the quota made sugar dollars certain? Cuban
manufacturing outside sugar processing, tobacco and a handful of consumer goods
was thin, and what existed was largely assembly.

It froze the land. Sugar wants large flat holdings and cheap seasonal labour,
and it got both. By the 1946 agricultural census something like 8 per cent of
farms held about 70 per cent of the farmland, and a substantial share of the
cane land --- the standard estimates run from a quarter to two-fifths of milling
capacity --- was owned by American companies.\unv{} Much of the best land was
held out of production entirely as a reserve against a good year, in a country
that imported a third of its food.

And it made unemployment structural. The cane harvest, the \emph{zafra}, ran
roughly January to May. The rest of the year was the \emph{tiempo muerto}, the
dead season, in which several hundred thousand cane workers had no work at all.
Annual average unemployment in the mid-1950s ran around 16 per cent and the
seasonal swing was far wider than that.\unv{} A country cannot build a stable
politics on an economy that lays off a fifth of its workforce every June.

To this was added the ownership question, which mattered less economically than
it did politically. American capital owned most of the electric and telephone
utilities, most of the nickel, a large share of the banking and the refineries,
and the milling capacity noted above. Whether this was extraction or investment
is a genuinely arguable question --- the utilities were real utilities and the
mills were real mills. What is not arguable is that it made the commanding
heights of the Cuban economy answerable to decisions taken in another country,
and that this was intolerable to a generation of Cubans who had grown up on the
nationalism of the 1933 revolution and the constitution of 1940.

\section{The republic that could not hold}

Cuba's 1940 constitution was, on paper, among the most advanced documents of
its era anywhere: it guaranteed the eight-hour day, paid holidays, the right to
strike, equal pay for women, free compulsory education, and it prohibited
latifundia and set limits on foreign landholding. Very little of the social
programme was ever implemented by enabling legislation, and the latifundia
clause was ignored outright. This gap between a superb constitution and a
government that did not execute it is the central fact of Cuban republican
politics, and it is why the constitution of 1940 became the banner of every
subsequent opposition, including Fidel Castro's.

The governments of Ram\'on Grau San Mart\'in (1944--48) and Carlos Pr\'io
Socarr\'as (1948--52) were elected, were not dictatorships, and were
spectacularly corrupt. Public money disappeared on a scale that was openly
discussed in the press and never prosecuted. Political violence between armed
factions, the \emph{grupos de acci\'on}, was routine in Havana. The Ortodoxo
party of Eduardo Chib\'as built a mass following on a single-issue platform of
public honesty --- his radio slogan was \emph{verg\"uenza contra dinero},
shame against money --- and when Chib\'as shot himself on air in August 1951 the
country lost the one figure who might have channelled the anti-corruption
current into an electoral outcome. Castro was an Ortodoxo candidate for
congress in the election that was never held.

On 10 March 1952, three months before that election, Fulgencio Batista, who had
run Cuba from behind and then in front between 1933 and 1944 and who was
polling third, took power in a bloodless dawn coup. This is the hinge. The coup
destroyed the constitutional route, and it destroyed it against a president who
was corrupt but elected and a constitution the country was proud of. Every
insurrectionary argument made in Cuba between 1952 and 1959 begins there, and
the argument was correct: after March 1952 there was no legal way to remove the
government.

Batista's second period was a dictatorship, and it got worse as it went on.
Congress was suspended, then restored as a rubber stamp through a rigged 1954
election that the opposition boycotted; the press was censored in waves; and
after the guerrilla insurgency began the police and army responded with
torture, disappearance and public killings, most notoriously by the Havana
police under Esteban Ventura and by the army in Oriente. The number of people
killed by the regime between 1952 and 1958 is genuinely unknown. The figure of
20,000, which entered circulation via a Havana newspaper in January 1959 and
was then repeated for sixty years in official Cuban texts, has never been
supported by any count and is not believed by serious historians in any camp;
careful estimates run from about two thousand to about five thousand, including
combatants \citep{thomas1971}.\unv{} That is a smaller number and an entirely sufficient one.

Meanwhile Havana became the offshore pleasure district of the United States. It
had had casinos before; under Batista, and under an explicit 1955 law offering
gaming licences and tax holidays to hotel builders, American organised crime ---
Meyer Lansky and Santo Trafficante most prominently --- became a structural
part of the tourist economy, with the regime taking its share. This is where the
tourism chapter of Part III has to start, because the memory of that Havana is
still doing political work in Cuba seventy years later, and any proposal to make
tourism the country's anchor industry has to answer it.

\section{Race, which the republic did not solve}

The 1953 census recorded roughly 27 per cent of Cubans as black or
mixed-race by self-identification, a figure most demographers regard as an
undercount.\unv{} Slavery had ended in 1886, later than anywhere in the Americas
except Brazil, and within living memory of people alive in 1958. The republic
had abolished legal racial distinction and had produced a genuinely
multiracial army, police and political class --- Batista himself was of mixed
race and was blackballed from Havana's white social clubs while running the
country.

What it had not abolished was the private colour bar: the beaches, clubs,
hotels and neighbourhoods that were white by practice; the hiring in banking,
commerce and the professions that was white by preference; the shape of the job
ladder in which Afro-Cubans were concentrated in cane cutting, construction,
domestic service and the informal economy. An attempt to legislate against
discrimination in employment in the late 1940s failed \citep{delafuente2001}. When the revolution
desegregated the beaches and the clubs in 1959 it was addressing something
real, it did so quickly, and it earned a loyalty among Afro-Cubans that has
outlasted almost every other source of the government's support --- a fact that
Part III has to take seriously, because the constituency most exposed to a
badly designed transition is the same one.

\section{The insurrection}

The armed opposition to Batista was not one organisation and Castro's was not
initially the largest.

It began on 26 July 1953 with an attack by about 130 young Ortodoxos on the
Moncada barracks in Santiago, which failed completely; roughly half the
attackers were killed, most after capture. Castro's trial defence, later
published as \emph{La historia me absolver\'a}, laid out a programme --- restore
the 1940 constitution, agrarian reform, profit-sharing, education, nationalise
the utilities --- that was radical-nationalist and not communist, and that
Cubans could and did read as a continuation of Chib\'as. He was sentenced,
amnestied in 1955 in a gesture of confidence by a regime that did not expect
what followed, and went to Mexico.

The \emph{Granma} landed 82 men in Oriente on 2 December 1956; a dozen or so
reached the Sierra Maestra. From there the July 26th Movement fought a small
guerrilla war --- never more than a few hundred fighters in the mountains ---
while a much larger urban underground, the \emph{llano}, ran sabotage, finance,
propaganda and supply in the cities under Frank Pa\'is until his killing by
police in July 1957 \citep{sweig2002}. A separate organisation, the Directorio Revolucionario of
the university students, nearly killed Batista in an assault on the
Presidential Palace in March 1957 and lost most of its leadership doing it. The
old Communist party, the Partido Socialista Popular, opposed the armed struggle
until late and joined it only in 1958, a fact that would matter enormously in
the 1960s.

Two things decided the war, and neither was military in the ordinary sense.
The first was that the general strike the movement called in April 1958 failed,
which pushed the centre of gravity decisively from the cities to the Sierra and
from the \emph{llano}'s leadership to Castro's. The second was that in March
1958 the United States, under congressional pressure over the regime's
atrocities, embargoed arms sales to Batista. The army was not defeated in the
field so much as it stopped being an army: the summer 1958 offensive into the
Sierra collapsed, whole units surrendered or changed sides, and when Che
Guevara's column took Santa Clara at the end of December the road to Havana was
open. Batista flew to the Dominican Republic in the small hours of 1 January
1959.

\section{What was actually inherited}

It is worth stating plainly, because both the later histories distort it.

\begin{ledger}{the inheritance, 1 January 1959}
\emph{Assets.} A middle-income economy with functioning infrastructure, a
literate urban majority, a real professional class, strong unions, a dense rail
and road network, a substantial sugar industry with modern mills, mineral
resources, and a tourism plant that already worked. A constitution with mass
legitimacy that everyone claimed to be restoring. Low external debt.

\emph{Liabilities.} A monoculture export economy dependent on a single
purchaser's political goodwill; structural seasonal unemployment near a fifth
of the workforce; a rural half of the country at Haitian levels of nutrition,
literacy and housing; concentrated landownership with much of the best land
idle; foreign ownership of the utilities, mines and much of the milling; an
unsolved colour bar; a political culture in which corruption was assumed and
the constitutional route had just been closed by force; and no institution
whatever with the legitimacy to arbitrate what came next --- the army was
Batista's, the congress was a fiction, the courts were compromised, and the
parties had been irrelevant since 1952.
\end{ledger}

That last liability is the one that decided the shape of everything after. In
January 1959 the only body in Cuba with unambiguous legitimacy was the
guerrilla movement, and within it, one man. There was no counterweight, and
there had been no counterweight since 10 March 1952. A great deal of what
follows in this book --- the speed of the radicalisation, the absence of any
internal brake on it, the ease with which a plural revolutionary coalition
became a single party --- is downstream of the fact that Batista's coup had
already destroyed the institutions that might have slowed it down. The
revolution did not seize a constitutional order. It inherited a vacuum where
one had been.

This is the first lesson the book will carry into Part III, and it is not an
abstract one: \emph{the sequence in which institutions are rebuilt matters more
than the speed}, and a transition that produces one legitimate actor and no
others will end where 1959 ended, whatever its intentions. Cuba has now run
that experiment twice, in 1952 and in 1959, from opposite directions.
